\documentclass[11pt]{article}

\usepackage[margin=1in]{geometry}

\begin{document}

\section*{New Completed PromptEvo Experiments}

\paragraph{Reporting scope.}
This note includes PromptEvo runs completed after the previous LaTeX record,
where Base, Stage 1, and Stage 2 were all finished on the same evaluation set.
The API-Bank rows are complete benchmark runs. The GEPA-AIME row is a completed
train-split diagnostic run, not a full AIME benchmark result. Ongoing tau2-bench
runs and baseline-only runs are reported as progress updates, not as completed
evolution results.

\begin{table}[htbp]
\centering
\small
\caption{Completed PromptEvo runs for the weekly report}
\label{tab:promptevo-weekly-completed-20260831}
\begin{tabular}{lrrr}
\hline
Experiment & Base & Stage 1 & Stage 2 \\
\hline
API-Bank protocol patch & 82.26\% & 83.29\% & 83.80\% \\
API-Bank tool-search subset & 54.62\% & 53.78\% & 56.30\% \\
GEPA-AIME & 20.00\% & 15.56\% & 20.00\% \\
\hline
\end{tabular}
\end{table}

\paragraph{API-Bank protocol patch LongCat 20260814.}
This run is complete, with all three stages evaluated on the same evaluation
set. Success rate improved from 82.26\% at Base to 83.29\% at Stage 1 and
83.80\% at Stage 2. Tool-call F1 improved from 0.939 to 0.960, parse success
improved from 97.17\% to 98.71\%, and runtime error rate was 0. This is the
strongest completed positive PromptEvo result for the weekly report.

\paragraph{API-Bank tool-search protocol patch LongCat 20260815.}
This run is complete, with all three stages evaluated on the same evaluation
set. Success rate changed from 54.62\% at Base to 53.78\% at Stage 1, then
recovered to 56.30\% at Stage 2. Tool-call F1 changed from 0.907 at Base to
0.898 at Stage 2, and parse success declined from 97.48\% to 95.80\%. Runtime
error rate was 0. This should be reported as a completed diagnostic subset: the
final success rate is above Base, but tool-call quality did not improve.

\paragraph{GEPA-AIME LongCat 20260815.}
This run is complete on the adapter's training split diagnostic setting, with
all three stages evaluated on the same evaluation set. Accuracy changed from
20.00\% at Base to 15.56\% at Stage 1, then returned to 20.00\% at Stage 2.
Format error, provider error, and runtime error rates were all 0. This result
should be treated as a small diagnostic run, not as a full AIME benchmark result
or evidence of positive improvement.

\paragraph{Suggested weekly-report wording.}
We completed three additional PromptEvo/LongCat experiments. The strongest
positive result is the full API-Bank protocol-patch run: success rate improved
from 82.26\% to 83.80\%, tool-call F1 improved from 0.939 to 0.960,
and parse success improved from 97.17\% to 98.71\%. The API-Bank tool-search
subset was also completed and showed mixed behavior: Stage 1 was below Base,
while Stage 2 recovered to 56.30\% and slightly exceeded Base. The GEPA-AIME
LongCat diagnostic run was completed with no provider or runtime errors, but
Stage 2 only returned to the Base accuracy of 20.00\%.

\paragraph{Current running item not counted as complete.}
The tau2-bench LongCat 20260816 run should not yet be reported as a completed
PromptEvo experiment. Its Base stage is complete, but Stage 1 is still being
resumed, and Stage 2/Stage 3 are not complete. In the weekly report, it should
only be described as: tau2-bench LongCat baseline completed; Stage 1 is still
running.

\end{document}
